% !TEX root =  master.tex
\chapter{Einleitung}

%\nocite{*}

Der Dartsport hat sich von einer einfachen Dartscheibe, bei welcher die Spielstände manuell erfasst werden mussten, schrittweise zu elektronischen Dartgeräten weiterentwickelt. Bei diesen Geräten wird der Spielbetrieb durch technische Funktionen unterstützt, 
etwa durch die automatische Erkennung von Treffern und das selbstständige Zählen der Punkte. Dadurch konnten klassische Spielabläufe vereinfacht und unterschiedliche Spielmodi direkt am Gerät bereitgestellt werden. Mit der zunehmenden Vernetzung der Geräte entsteht darüber hinaus ein erweitertes System, 
bei dem das Dartgerät nicht mehr isoliert betrieben wird, sondern mit weiteren digitalen Komponenten verbunden ist. 

Das E-Dart-Gerät \textit{LÖWEN DART HB10} (s. \cref{fig:hb10}) stellt in diesem Zusammenhang eine aktuelle Weiterentwicklung dar. LÖWEN DART prägt seit 1986 den elektronischen Dartsport 
und verbindet mit dem HB10 traditionelle E-Dart-Funktionen mit digitalen Erweiterungen. Gerade im Breitensport ist Dart heute sehr beliebt. Dadurch haben sich über den Verlauf der letzten Jahre nationale Ligen und Turnier-Serien etabliert, die durch offizielle Verbände, 
wie den \ac{DSAB}\footnote{https://dsab-vfs.de/} veranstaltet werden. Der HB10 wird als offizielles \ac{DSAB}-Liga- und Turniergerät eingesetzt und bietet verschiedene Spielvarianten, Trainingsspiele sowie einen Turniermodus. 
Zu den technischen und funktionalen Eigenschaften zählen unter anderem ein hochauflösender Monitor, ein Touchpad, ein Treffererkennungssystem, eine App-Anmeldung über einen QR-Code im Spielermenü sowie die Möglichkeit zur Netzwerkanbindung \cite{lwenentertainment_2026_lwen}.

\begin{figure}[h]
    \centering
    \includegraphics[height=0.3\textheight]{Images/HB10front.png}
    \caption[HB10-Dartgerät]{Das LÖWEN DART HB10 dient im entwickelten System als Spielgerät und Ausgangspunkt der Videoaufzeichnung. Das Gerät erfasst die während eines Spiels auftretenden Treffer und wird im Rahmen dieser Arbeit um eine Kamera sowie eine Anwendung zur Verarbeitung der Videodaten erweitert.}
    \label{fig:hb10}
\end{figure}

Durch diese Vernetzung entsteht ein System, 
das über die reine Funktion einer elektronischen Dartscheibe hinausgeht. Während ein Offline-E-Dart-Gerät Spiele bereitstellt und Punkte automatisch zählt, ermöglicht ein vollständig vernetztes Gerät die Einbindung in einen größeren digitalen Kontext, 
bestehend aus Gerät, Server und App. In diesem Umfeld können spielbezogene Daten zentral verarbeitet und für weitere Funktionen genutzt werden. Für Spieler ergeben sich daraus unter anderem Statistiken, Ranglisten und vernetzte Turnierformate.  
Zusätzlich steht für Aufsteller mit dem LÖWEN DART-Portal ein zentraler Online-Zugriff auf HB10-Geräte zur Verfügung, über den Umsätze, Gerätedaten und Statistiken eingesehen werden können \cite{lwenentertainment_2026_lwen}. 

Ein zentraler Aspekt der Updates der letzten beiden Jahre des HB10 sind vernetzte Turniere, bei denen Spielergebnisse digital verarbeitet und im Anschluss an ein Spiel an einen Server übertragen werden. 
Diese werden in Kooperation mit 3K Dart \cite{mützner_2026} komfortabel online organisiert und dann vor Ort an HB10 ausgespielt. Über diese Plattform können sowohl Turniere als auch Ligen ausgespielt werden. 
Das größte Turnier dieser Art fand vor kurzem mit über 4.000 aktivem Spielern und über 250 HB10 statt \cite{2026}. Diese synchronen Turnierformen setzen voraus, dass die Teilnehmer innerhalb eines gemeinsamen zeitlichen Rahmens und am selben HB10 einzelne Runden des Wettbewerbs austragen. 
In Zukunft soll dieses System um asynchrone Turnierformate und Ranglisten ergänzt werden, bei denen die Möglichkeit besteht, Preisgelder zu gewinnen. Im Gegensatz zu den bestehenden Formaten ist hierbei keine gleichzeitige Anwesenheit der Spieler erforderlich. 
Da die Wettbewerbe ortsungebunden durchgeführt werden können, ist eine spätere visuelle Überprüfung der sportlichen Integrität erforderlich. 

Ein mögliches asynchrones Spiel kann dabei so verstanden werden, dass ein Spieler ein Spiel eigenständig an einem HB10 durchführt, ohne dass andere Teilnehmer zeitgleich anwesend sein müssen. Während des Spielverlaufs werden die einzelnen Würfe 
durch in das Gerät integrierte Kameras erfasst und als kurze Videosequenzen zwischengespeichert. Nach Abschluss des Spiels werden diese einzelnen Videoaufnahmen in eine zeitlich geordnete Form gebracht, zu einem zusammenhängenden Video zusammengefügt und anschließend an einen Server übertragen.
Dort wird das Video gespeichert und so bereitgestellt, dass andere Spieler oder berechtigte Nutzer es zu einem späteren Zeitpunkt über eine Webplattform abrufen und ansehen können. Dadurch entsteht ein Spielablauf, bei dem die sportliche Leistung nicht nur über das gemeldete Spielergebnis, 
sondern zusätzlich über eine nachträglich überprüfbare visuelle Dokumentation nachvollzogen werden kann. 

Mit dieser Erweiterung ergeben sich zusätzliche technische und organisatorische Anforderungen an das Gesamtsystem. Auf Seiten des HB10 muss eine geeignete Kamera integriert werden, die während des Spielbetriebs verwertbare Bilddaten liefert, ohne den eigentlichen Spielablauf zu beeinträchtigen. 
Darüber hinaus müssen die aufgenommenen Videodaten zwischengespeichert, verarbeitet, komprimiert und nach Spielende zuverlässig übertragen werden. Auf Serverseite ist eine Speicherung der Videos sowie der zugehörigen Metadaten erforderlich, damit die Aufnahmen eindeutig einem Spiel, 
einem Zeitpunkt und gegebenenfalls beteiligten Spielern zugeordnet werden können. Zusätzlich muss der Zugriff auf die Videos über eine Webplattform geregelt werden, sodass nur vorgesehene Nutzer die gespeicherten Inhalte abrufen können. 
Da bei der Aufnahme und Bereitstellung von Videomaterial personenbezogene Informationen entstehen können, sind außerdem Anforderungen an Datenschutz, Zugriffskontrolle und den Umgang mit gespeicherten Aufnahmen zu berücksichtigen.

Diese Arbeit befasst sich mit dem Konzept und der Umsetzung eines Prototypens für eine \ac{VOD}-Streaming-Plattform für das \textit{LÖWEN DART HB10}, 
die eine Überprüfung vergangener Spiele durch die Nutzer ermöglicht. Um dies zu erreichen, wird zunächst in \cref{sec:system} ein Konzept für das Streaming-System entwickelt. Daraufhin wird in \cref{sec:dart} beschrieben, wie eine Kamera in den HB10 integriert und das Gerät um die Funktion erweitert wurde, 
die von der Kamera aufgenommenen Videos für das \ac{VOD}-Streaming aufzubereiten. Anschließend wurde in \cref{sec:server} ein Server entwickelt, welcher die Videos des HB10 empfängt, zwischenspeichert und auf Anfrage des Nutzers bereitstellt. In \cref{sec:web} wird Beschrieben, 
wie eine Webplattform erstellt wird, über welche die Nutzer die gespeicherten Videos ansehen können.

\newpage
